% Source-derived chapter 04: Line bundles over abelian schemes
% Original source: geometric-bogomolov-conjecture-arbitrary-characteristics
\chapter{Line bundles over abelian schemes}
\label{geometric-bogomolov-conjecture-arbitrary-characteristics-chapter-04}
\source{geometric-bogomolov-conjecture-arbitrary-characteristics}

\begin{remark}[Source section]
\label{geometric-bogomolov-conjecture-arbitrary-characteristics-chapter-04-source}
\source{geometric-bogomolov-conjecture-arbitrary-characteristics}
\source{geometric-bogomolov-conjecture-arbitrary-characteristics}
This chapter reproduces the corresponding section of the retrieved
LaTeX source, including its definitions, statements, examples, and
proofs. It has not yet been translated into Lean.
\notready
\end{remark}

% The source section title was: Line bundles over abelian schemes
\label{sec line bundle}
\source{geometric-bogomolov-conjecture-arbitrary-characteristics}

In this section, we introduce some preliminary results for abelian schemes over curve, which will be used in the proof the geometric Bogomolov conjecture in \S\ref{sec final}. These results are well-known, but we collect them here for convenience.

\subsection{Rigidified line bundles}
Let $S$ be a smooth projective curve over a field $k$. 
Let $\pi:\sA\to S$ be an abelian scheme over $S$.
This will be the basic setup of this section.

For any $m\in\Z$, denote by $[m]:\sA\to \sA$ the homomorphism of multiplication by $m$, and denote by $\sA[m]$ the kernel of this homomorphism. 

By a \emph{multi-section} of $\pi:\sA\to S$, we mean a closed integral subscheme $\sT$ of $\sA$ such that the induced morphism $\sT\to S$ is finite and flat. 
The multi-section $\sT$ is called \emph{torsion} if its corresponding element in the abelian group $\sA_\sT(\sT)$ is torsion. 
The \emph{order} of the torsion multi-section $\sT$ is defined to be the order of the corresponding element in $\sA_\sT(\sT)$.

Any torsion multi-section $\sT$ is necessarily the Zariski closure of a torsion point of $A(\OK)_\tor$ in $\sA$. 
If the order of $\sT$ is $m$, then $\sT$ is a closed subscheme of $\sA[m]$. 
If furthermore $m$ is not divisible by $\mathrm{char}(k)$, then $\sA[m]$ is finite and \'etale over $S$, and thus $\sT$ is also finite and \'etale over $S$. 
In this case, $\sT$ is actually smooth over $k$.

Let $\sL$ be a line bundle over $\sA$. 
We say that $\sL$ is \emph{symmetric} if $[-1]^*\sL\simeq \sL$. 
We say that $\sL$ is \emph{anti-symmetric} if $[-1]^*\sL\simeq \sL^\vee$. 
We say that $\sL$ is \emph{rigidified} if it is endowed with an isomorphism $e^*\sL\simeq \sO_S$ via the identity section $e:S\to \sA$. The isomorphism is called a \emph{rigidification}.
The following results are well-known to experts, but we sketch proofs for convenience of readers.

\begin{lemma}
\source{geometric-bogomolov-conjecture-arbitrary-characteristics}
\notready \label{line bundle}
\source{geometric-bogomolov-conjecture-arbitrary-characteristics}
Let $\sL$ be a rigidified line bundle over $\sA$. 
Then the following hold.
\begin{itemize}
\item[(1)] If $\sL$ is symmetric, then $[m]^*\sL\simeq m^2\sL$ for any $m\in \Z$; if $\sL$ is anti-symmetric, then $[m]^*\sL\simeq m\sL$ for any $m\in \Z$.
\item[(2)] For any torsion multi-section $\sT\subset \sA$, the line bundle $\sL|_{\sT}$ is torsion in $\Pic(\sT)$.
\item[(3)] If $\sL$ is symmetric and $\pi$-ample, then $\sL$ is nef over $\sA$. 
\item[(4)] If $\sL$ is symmetric and $\pi$-ample, and $\sL'$ is another symmetric and rigidified line bundle over $\sA$, then there is a positive integer $a$ such that $a\sL-\sL'$ is nef over $\sA$. 
\end{itemize}
\end{lemma}

\begin{proof}
For (1), we only consider the symmetric case as the anti-symmetric case is similar. 
Note that $[m]^*\sL- m^2\sL$ is trivial over every fiber of $\pi:\sA\to S$, so it lies in $\pi^*\Pic(S)$. Consider the pull-back to the identity section. Then $[m]^*\sL- m^2\sL$ is trivial by the rigidification. 

For (2), by $2\sL=(\sL+[-1]^*\sL)+(\sL-[-1]^*\sL)$, we can assume that $\sL$ is either symmetric case or anti-symmetric. Let $m$ be the order of $\sT$. It suffices to prove that $\sL|_{\sA[m]}$ is torsion. Note that $\sA[m]\to S$ is the base change of $[m]:\sA\to \sA$ by the identity section. 
It follows that $([m]^*\sL)|_{\sA[m]}$ is trivial over $\sA[m]$ by the rigidification. 
Then $\sL|_{\sA[m]}$ is torsion over $\sA[m]$ by (1).

For (3), since $\sL$ is $\pi$-ample, there is an ample line bundle $\sL'$ over $S$ such that $\sL+\pi^*\sL'$ is ample over $\sA$. 
Then $[m]^*(\sL+\pi^*\sL')\simeq m^2\sL+\pi^*\sL'$ is ample over $\sA$.
The $\Q$-line bundle $\sL+m^{-2}\pi^*\sL'$ is ample over $\sA$, and thus its limit $\sL$ is nef over $\sA$.

For (4), note that $a\sL-\sL'$ is $\pi$-ample for sufficiently large $a$. Apply (3). 
\end{proof}


\subsection{Numerical classes of torsion multi-sections}

The following actually holds for rational equivalence and for high dimensional base $S$ (cf. \cite[Theorem 2.1]{GZ2014}). 
We will be content to numerical equivalence for $\dim S=1$, which has the following quick proof and is sufficient for our application.

\begin{pro} \label{numerical equivalent}
\source{geometric-bogomolov-conjecture-arbitrary-characteristics}
Let $S$ be a smooth projective curve over a field $k$. 
Let $\pi:\sA\to S$ be an abelian scheme over $S$ of relative dimension $g\geq 1$.
Let $\sL$ be a symmetric and rigidified line bundle over $\sA$. 
Let $\sT$ be a torsion multi-section of $\sA$ over $S$. 
Then there is a numerical equivalence 
$$[\sL]^g\equiv
\frac{\deg(\sL_\eta)}{\deg(\sT/S)} [\sT]$$ 
of 1-cycles over $\sA$. 
Here $\eta$ is the generic point of $S$, and $\deg(\sL_\eta)$ is the degree of $\sL_\eta$ over the abelian variety $\sA_\eta$. 
\end{pro}

\begin{proof}
It suffices to prove that 
$$\sL^g\cdot \sM=
\frac{\deg(\sL_\eta)}{\deg(\sT/S)} \sT\cdot \sM$$ 
for any line bundle $\sM$ over $\sA$. 

We first prove that if $\sM$ is rigidified, then both sides are 0. 
The right-hand side is 0 by Lemma \ref{line bundle}(2).
For the left-hand side, by the decomposition $2\sL=(\sL+[-1]^*\sL)+(\sL-[-1]^*\sL)$ again, we can assume that $\sL$ is either symmetric case or anti-symmetric. 
Then $[m]^*\sM\simeq m^i\sM$ for $i=1,2$.
By the projection formula, 
$$([m]^*\sL)^g\cdot ([m]^*\sM)= 
\deg([m]) \sL^g\cdot \sM.$$ 
This is just 
$$m^{2g+i} \,\sL^g\cdot \sM= 
m^{2g}\, \sL^g\cdot \sM.$$ 
It follows that $\sL^g\cdot \sM=0.$

To prove the identity for general $\sM$, write $\sM=(\sM-\pi^*e^*\sM)+\pi^*e^*\sM$. Note that $\sM-\pi^*e^*\sM$ is canonically rigidified. It is reduced to prove the result for $\sM=\pi^*\sN$ for line bundles $\sN$
over $S$. This follows from the simply equalities
$$\sL^g\cdot \pi^*\sN=\deg(\sL_\eta)\deg(\sN),\quad
\sT\cdot \pi^*\sN=\deg(\sT/S)\deg(\sN).$$ 
This finishes the proof.
\end{proof}



\subsection{Canonical height}

Let $S$ be a smooth projective curve over a field $k$. 
Let $K=k(S)$ be the function field. 
Let $A$ be an abelian variety over $K$, and $L$ is a symmetric and ample line bundle over $A$.
Let $X$ be a closed subvariety of $A$. 
Then we have the canonical height $\hat h(X)$ associated to $L$. 

Assume that $A$ has everywhere \emph{good reduction} over $S$.
In this case, we have the following well-known easy interpretation of the canonical height of closed subvarieties of $A$. 

In fact, let $\pi:\sA\to S$ be the unique abelian scheme with generic fiber $A\to \Spec K$.
Let $\sL$ be a symmetric and rigidified line bundle over $\sA$ extending $L$.

For the existence of $\sL$, we first take any line bundle $\sL'$ over $\sA$ extending $L$, which can be obtained by passing to divisors and taking Zariski closures. 
By replacing $\sL'$ by $\sL'-\pi^*e^*\sL'$, we can assume that $\sL'$ is rigidified.  Here $e:S\to\sA$ is the identity section.
Then $\sL$ is automatically symmetric. 
In fact, $\sL-[-1]^*\sL$ is trivial over $A$, so it is isomorphic to $\pi^*\sM$ for some line bundle $\sM$ over $S$.
The rigidification implies that $\sM$ is trivial. 

Denote by $\sX$ the Zariski closure of $X$ in $A$. 
By Lemma \ref{line bundle}(1) and the projection formula, the naive height
$$
h_{\sL}(X)=\frac{\sL^{\dim X+1}\cdot \sX}{(\dim X+1)\deg_L(X)}
$$
already satisfies
$h_{\sL}([n]X)=n^2\, h_{\sL}(X)$. 
As a consequence, $h_{\sL}(X)$ is exactly equal to the canonical height 
$\hat h(X)$ associated to $L$.
